%! AP Statistics — Unit 7, Lesson 7.2 — Exit Ticket (Answer Key)
\documentclass[10pt]{article}
\usepackage{apstats-article}
\usepackage{apstats-key}
\newcommand{\TallMath}[1]{$\displaystyle #1\rule[-1.4em]{0pt}{3.2em}$}

\begin{document}

\pageheader{Unit 7, Lesson 7.2}{Exit Ticket --- Answer Key}
\namedateperiod

\vspace{0.1in}
A quality-control engineer randomly selects 16 bottles from a production line and
measures the fill volume. $\bar x = 498$ mL, $s = 6$ mL.

\begin{enumerate}

  \item Conditions:
        \begin{itemize}
            \item \textbf{Random:} \ans{Random sample of bottles stated.}
            \item \textbf{10\%:} \ans{$16 \le 10\%$ of all bottles produced (assume large production run).}
            \item \textbf{Large Sample:} \ans{$n = 16 < 30$; must check for strong skewness or outliers in the data. Assume engineer verified this.}
        \end{itemize}

  \item $df = $ \ans{15} \quad $t^* = $ \ans{2.131} \quad
        $SE = 6/\sqrt{16} = $ \ans{1.5 mL}

  \item $498 \pm 2.131 \times 1.5 = 498 \pm 3.197 = ($ \ans{494.8} $,$ \ans{501.2} $)$ mL

  \item \ansline{We are 95\% confident that the interval (494.8, 501.2) mL captures the true mean fill volume for all bottles produced on this production line.}

\end{enumerate}

\begin{teachernote}
    $t^*$ for $df = 15$, 95\% CI is 2.131.
    $ME = 2.131 \times 1.5 = 3.197$ mL.
\end{teachernote}

\end{document}
